% Einheit 2 von 5 – Quader und Würfel (bank/koerper/e2.jsonl, zusammenbau v0.1)
%% TODO Zweigzeile Teil 1: was hier gelernt wird (Satz in Schülersprache aus dem Einheitstitel) – Einheit 2
\einheitenkopf[e2]{Einheit 2 von 5 · Quader und Würfel}
\zweigzeile{ab Kl. 5, je nach Buch bis Kl. 6 · P10 oft}
\setcounter{aufgabe}{13}

%% TODO Titel: Ich-kann-Satz fehlt in der Bank (Platzhalter: Kettenname) – Nr. 14
%% TODO Anweisung: ein Satz über der ersten Teilaufgabe fehlt in der Bank – Nr. 14
\begin{aufgabe}{Quader}
\begin{teile}
\teil Wie viel Wasser passt in das Aquarium? Gefragt ist das … \kreuz{Volumen} \kreuz{Oberfläche}
\end{teile}
%% TODO Darstellung für schwach fehlt in der Bank (koerper-e2-k1-s1-v1; Abschnitt „Für schwache Schüler“)
\swz{Quader: Länge $6$ cm, Breite $3$ cm, Höhe $2$ cm – Volumen?}{}
%% TODO Darstellung für schwach fehlt in der Bank (koerper-e2-k1-s1-v2; Abschnitt „Für schwache Schüler“)
\swz{Quader: Länge $7$ cm, Breite $5$ cm, Höhe $2$ cm – Volumen?}{}
%% TODO Darstellung für schwach fehlt in der Bank (koerper-e2-k1-s1-v3; Abschnitt „Für schwache Schüler“)
\swz{Container: Länge $8$ m, Breite $5$ m, Höhe $3$ m – Volumen?}{}
%% TODO Darstellung für schwach fehlt in der Bank (koerper-e2-k1-s1-v4; Abschnitt „Für schwache Schüler“)
\swz{Kiste: Länge $9$ dm, Breite $4$ dm, Höhe $3$ dm – Volumen?}{}
%% TODO Darstellung für schwach fehlt in der Bank (koerper-e2-k1-s1-v5; Abschnitt „Für schwache Schüler“)
\swz{Quader: Länge $5$ cm, Breite $5$ cm, Höhe $3$ cm – Volumen?}{}
\swfrage{Was bleibt gleich, was ändert sich?}
%% TODO Darstellung für schwach fehlt in der Bank (koerper-e2-k1-s2-v1; Abschnitt „Für schwache Schüler“)
\swz{Ein Würfel hat $5$ cm Kantenlänge – Volumen?}{}
%% TODO Darstellung für schwach fehlt in der Bank (koerper-e2-k1-s3-v1; Abschnitt „Für schwache Schüler“)
\swz{$4\,500$ cm³ – wie viele Liter?}{}
%% TODO Darstellung für schwach fehlt in der Bank (koerper-e2-k1-s4-v1; Abschnitt „Für schwache Schüler“)
\swz{Quader: $7$ cm lang, $3$ cm breit, $2$ cm hoch – Oberfläche?}{}
%% TODO Darstellung für schwach fehlt in der Bank (koerper-e2-k1-s5-v1; Abschnitt „Für schwache Schüler“)
\swz{Ein würfelförmiger Hocker hat $5$ dm Kantenlänge – Oberfläche?}{}
%% TODO Darstellung für schwach fehlt in der Bank (koerper-e2-k1-s6-v1; Abschnitt „Für schwache Schüler“)
\swz{Quader: Volumen $90$ cm³, Länge $5$ cm, Breite $3$ cm – Höhe?}{}
%% TODO Darstellung für schwach fehlt in der Bank (koerper-e2-k1-s7-v1; Abschnitt „Für schwache Schüler“)
\swz{Eine Betontreppe besteht aus zwei Quadern: unten $6$ dm × $4$ dm × $2$ dm, oben darauf $3$ dm × $4$ dm × $2$ dm. Volumen?}{}
\swz[3]{Ein Würfel hat die Kantenlänge $8$ cm. Gib sein Volumen an \hfill (P10 2026 FOR).}{}
\end{aufgabe}

%% TODO Titel: Ich-kann-Satz fehlt in der Bank (Platzhalter: Kettenname) – Nr. 15
%% TODO Anweisung: ein Satz über der ersten Teilaufgabe fehlt in der Bank – Nr. 15
\begin{aufgabe}{Oberfläche ohne Deckel (Kiste, Aquarium)}
%% TODO Darstellung für schwach fehlt in der Bank (koerper-e2-k2-s1-v1; Abschnitt „Für schwache Schüler“)
\swz{Eine Holzkiste ohne Deckel ist $8$ dm lang, $5$ dm breit und $4$ dm hoch. Wie viel Holz braucht man?}{}
\end{aufgabe}

%% TODO Titel: Ich-kann-Satz fehlt in der Bank (Platzhalter: Kettenname) – Nr. 16
%% TODO Anweisung: ein Satz über der ersten Teilaufgabe fehlt in der Bank – Nr. 16
\begin{aufgabe}{Würfelkante aus V (Kubikzahlen)}
%% TODO Darstellung für schwach fehlt in der Bank (koerper-e2-k3-s1-v1; Abschnitt „Für schwache Schüler“)
\swz{Würfel mit $27\,000$ cm³ Volumen – Kantenlänge?}{}
\end{aufgabe}

%% TODO Titel: Ich-kann-Satz fehlt in der Bank (Platzhalter: Kettenname) – Nr. 17
%% TODO Anweisung: ein Satz über der ersten Teilaufgabe fehlt in der Bank – Nr. 17
\begin{aufgabe}{Einheiten cm³, dm³ = l, m³ (1 l = 1000 cm³)}
%% TODO Darstellung für schwach fehlt in der Bank (koerper-e2-k4-s1-v1; Abschnitt „Für schwache Schüler“)
\swz{Ein Pool fasst $3$ m³ – wie viele Liter?}{}
\end{aufgabe}

%% TODO Titel: Ich-kann-Satz fehlt in der Bank (Platzhalter: Kettenname) – Nr. 18
%% TODO Anweisung: ein Satz über der ersten Teilaufgabe fehlt in der Bank – Nr. 18
\begin{aufgabe}{Aquarium: Füllmenge in Litern, Füllhöhe}
%% TODO Darstellung für schwach fehlt in der Bank (koerper-e2-k5-s1-v1; Abschnitt „Für schwache Schüler“)
\swz{Ein Aquarium ist innen $50$ cm lang und $25$ cm breit. Es werden $50$ l Wasser eingefüllt. Wie hoch steht das Wasser?}{}
\end{aufgabe}

%% TODO Titel: Ich-kann-Satz fehlt in der Bank (Platzhalter: Kettenname) – Nr. 19
%% TODO Anweisung: ein Satz über der ersten Teilaufgabe fehlt in der Bank – Nr. 19
\begin{aufgabe}{Quader – Fehler finden, Begründen, Anwendung, Darstellungswechsel}
\swz[3]{Ein Quader ist $7$ cm lang, $2$ cm breit und $5$ cm hoch. Nils berechnet die Oberfläche: \rechnung{7 \cdot 2 \cdot 5 &= 70 \text{ cm}^2} Wo ist der Fehler?}{}
\swfrage{Warum ist das Volumen eines Quaders Länge mal Breite mal Höhe? Erkläre mit Würfeln von $1$ cm Kantenlänge.}
\swz[3]{Ein Umzugskarton ist innen $62$ cm lang, $34$ cm breit und $36$ cm hoch. Im Katalog steht „ca. 75 Liter“. Stimmt die Angabe?}{}
\swa{Zeichne das Netz eines Quaders mit $7$ cm, $2$ cm und $1$ cm Kantenlänge und schreibe in jedes Rechteck seinen Flächeninhalt.}{\rechenplatz{5}}
\end{aufgabe}

\uebersichtskasten{Volumen: Länge mal Breite mal Höhe. Würfel: Kante mal Kante mal Kante. \\ a = 5 cm, b = 4 cm, c = 2 cm: V = 5 · 4 · 2 = 40 cm³ \quad Würfel a = 4 cm: V = 4 · 4 · 4 = 64 cm³ \\ Oberfläche: sechs Rechtecke, je zwei gleich. \quad O = 2 · (5 · 4 + 5 · 2 + 4 · 2) = 76 cm² \quad Würfel: O = 6 · 4 · 4 = 96 cm² \\ Einheiten: 1 dm³ = 1 l = 1000 cm³ \quad 1 m³ = 1000 l \\ Rückwärts: c = V : (a · b).}
